\chapter{Neural Edge-Node Fusion: A Dual-Pipeline Architecture for Causal Discovery}
\label{chap:method}

This chapter describes the complete pipeline of the model proposed in this thesis. The method takes as input an observational dataset of $n = 1000$ observations over $p \in \{3,\ldots,10\}$ variables -- with a designated treatment variable $X$ and outcome variable $Y$ -- and produces an 8-class causal role label for each remaining variable $K \notin \{X, Y\}$.

The proposed model builds on the edge-centric deep learning framework introduced by the first-place solution of the ADIA Lab Causal Discovery Challenge~\cite{adia_rank1_report, olivetti2025adialab}, and extends it with five targeted contributions: the first contribution expands the edge tensor from 3 to 8 channels through multi-bandwidth kernel regression and ANM residuals, capturing causal asymmetry at multiple scales with a gain of +1.89\%; the second contribution injects topology-aware priors into self-attention layers using only six learned scalars, encoding chain, fork, and collider relationships directly to achieve +2.65\%; the third contribution introduces a 2D scatter density pipeline that renders joint distributions as visual images, capturing patterns invisible to pairwise analysis with +0.88\%; the fourth contribution uses a $\lambda=0.7$ dual-loss weighting scheme that ensures edge embeddings remain causally meaningful throughout training, adding +0.59\%; and the fifth contribution applies X/Y remap data augmentation that multiplies the training set by approximately 11x without additional data collection, contributing +1.55\% to +3.32\% depending on configuration.

% =====================================================================
\section{Architecture Overview}
\label{sec:overview}
% =====================================================================

The model consists of two parallel processing pipelines that operate on different representations of the same observational data, followed by a fusion step for final classification. Figure~\ref{fig:architecture} provides the complete architectural diagram illustrating how these components connect.

The \textbf{Edge Context Encoder} (Section~\ref{sec:edge_pipeline}) processes all $p(p-1)$ directed variable pairs simultaneously. Each ordered pair $(u, v)$ is represented as an 8-channel 1D tensor of length $n$, capturing the statistical and causal relationship between $u$ and $v$ through sorted observations, multi-bandwidth kernel regression curves, and ANM residuals. These tensors are processed independently by a shared-weight Conv1D network, then aggregated across the full graph via structural self-attention. A binary edge head supervises this pipeline during training, forcing the edge embeddings to encode causally meaningful graph structure.

The \textbf{Node Visual Classifier} (Section~\ref{sec:node_pipeline}) processes each non-$X$/$Y$ variable $K$ individually. For each such node, an 8-channel $32 \times 32$ scatter density image is constructed by rendering the joint distribution of $K$ against $X$ and $Y$ (and their ANM residuals) as smoothed 2D histograms. A hierarchical Conv2D network extracts visual features from this image.

The \textbf{Fusion} step (Section~\ref{sec:fusion}) combines the node's 2D visual embedding with the four edge context embeddings corresponding to edges $K \to X$, $K \to Y$, $X \to K$, and $Y \to K$, producing the final 8-class prediction for each node.

\input{figures/architecture_tikz}

% =====================================================================
\section{Preprocessing and Edge Tensor Construction}
\label{sec:preprocessing}
% =====================================================================

This section describes the construction of the 8-channel edge tensor representation, which is the first contribution of this thesis extending the 3-channel baseline.

\subsection{Data Normalization}

Each input dataset is a matrix $\mathbf{D} \in \mathbb{R}^{n \times p}$ with $n = 1000$ observations. Prior to feature construction, each column (variable) is standardized to zero mean and unit variance via z-score normalization:
\begin{equation}
\tilde{x}_i = \frac{x_i - \mu}{\sigma + \varepsilon},
\label{eq:zscore}
\end{equation}
where $\mu$ and $\sigma$ are the column mean and standard deviation respectively, and $\varepsilon = 10^{-8}$ prevents division by zero. Normalization ensures that all subsequent computations -- particularly kernel bandwidth selection -- are scale-invariant across variables and datasets.

\subsection{Complete Directed Graph and Edge Pairs}

From $p$ variables, a complete directed graph with $p(p-1)$ ordered pairs is constructed. Every ordered pair $(u, v)$ with $u \neq v$ produces a separate edge tensor: the pairs $(u, v)$ and $(v, u)$ are treated as distinct, because the directional asymmetry between them is a primary causal signal. This permutation-equivariant design ensures that classification results are independent of the arbitrary naming order of variables.

\subsection{8-Channel Edge Tensor}
\label{sec:8channel_tensor}

For each ordered pair $(u, v)$, let $\sigma$ denote the permutation that sorts variable $u$ in ascending order, so that $u_{\sigma(1)} \leq \cdots \leq u_{\sigma(n)}$. The 8-channel edge tensor $\mathbf{T}_{uv} \in \mathbb{R}^{8 \times n}$ is constructed as follows. Figure~\ref{fig:edge_tensor} illustrates the resulting curves for each channel.

\textbf{Channels 1--2: Sorted observations.} Channel 1 preserves the marginal distribution of $u$: $\mathbf{c}_1 = (u_{\sigma(1)},\ldots,u_{\sigma(n)})$. Channel 2 records $v$ reordered by the rank of $u$: $\mathbf{c}_2 = (v_{\sigma(1)},\ldots,v_{\sigma(n)})$, encoding the functional relationship between $u$ and $v$ as an empirical curve. When sorted by rank of $u$, channel 2 traces the conditional expectation $\mathbb{E}[v \mid u]$: monotone relationships appear as smooth increasing or decreasing curves, non-monotone relationships as curved or U-shaped traces. This representation retains the full shape of the conditional relationship, which is information irreversibly lost by scalar summaries such as correlation coefficients.

\textbf{Channels 3--5: Multi-bandwidth kernel regression.} At each sorted point $u_{\sigma(i)}$, the Nadaraya--Watson estimator~\cite{nadaraya1964, watson1964} computes:
\begin{equation}
\hat{m}_h\!\left(u_{\sigma(i)}\right) = \frac{\sum_{j=1}^{n}\exp\!\Bigl(-\frac{(u_{\sigma(i)}-u_j)^2}{2h^2}\Bigr)\,v_j}{\sum_{j=1}^{n}\exp\!\Bigl(-\frac{(u_{\sigma(i)}-u_j)^2}{2h^2}\Bigr)},
\label{eq:nadaraya}
\end{equation}
with three bandwidths $h \in \{0.2,0.5,1.0\}$, yielding channels $\mathbf{c}_3, \mathbf{c}_4, \mathbf{c}_5$. Each bandwidth captures a different smoothness scale of the regression function: $h = 0.2$ is sensitive to local curvature and sharp transitions, while $h = 1.0$ reflects global monotone trends. Using all three bandwidths together prevents the model from depending on a single fixed scale assumption, which is critical because the optimal smoothness varies across different causal mechanisms.

\textbf{Channels 6--8: Multi-bandwidth ANM residuals.} The Additive Noise Model~\cite{hoyer2009anm} posits that if $U \to V$ is the true causal direction, the residual $V - \hat{m}(U)$ should be approximately independent of $U$; the reverse-direction residual retains structure. This asymmetry is encoded in channels 6--8:
\begin{equation}
r_h\!\left(u_{\sigma(i)}\right) = v_{\sigma(i)} - \hat{m}_h\!\left(u_{\sigma(i)}\right), \quad h \in \{0.2,0.5,1.0\},
\label{eq:anm_resid}
\end{equation}
yielding channels $\mathbf{c}_6, \mathbf{c}_7, \mathbf{c}_8$. When $u \to v$ holds, the residual curve is approximately flat (independent of $u$); under the reversed direction, the residual retains a systematic nonlinear shape. The Conv1D network learns to detect this difference directly from the curve shape.

\input{figures/edge_tensor_tikz}

% =====================================================================
\section{Edge Context Encoder}
\label{sec:edge_pipeline}
% =====================================================================

This section describes the 1D pipeline that processes all edge tensors and aggregates them into node-level representations through structural self-attention.

\subsection{Per-Edge Feature Extraction via Conv1D}

The 8-channel edge tensor $\mathbf{T}_{uv} \in \mathbb{R}^{8 \times n}$ is processed by a shared-weight Conv1D network to produce a $d = 64$ dimensional edge embedding. The network consists of a stem layer followed by five residual blocks. The stem layer expands the 8 input channels to $d = 64$ hidden channels:
\begin{equation}
\mathbf{H}_0 = \mathrm{Conv1D}_{8 \to 64}(\mathbf{T}_{uv}) \in \mathbb{R}^{64 \times n}.
\label{eq:stem}
\end{equation}
Each of the five residual Conv1D blocks applies the pattern from He et al.\@\cite{he2016resnet}, adapted with Group Normalization~\cite{wu2018groupnorm} for stability under small batch sizes and the GELU activation~\cite{hendrycks2016gelu}:
\begin{equation}
\mathbf{H}_{l+1} = \mathbf{H}_l + \mathrm{Conv1D}\!\Bigl(\mathrm{GELU}\!\bigl(\mathrm{GroupNorm}(\mathbf{H}_l)\bigr)\Bigr).
\label{eq:resblock}
\end{equation}
After the five residual blocks, Global Average Pooling reduces the sequence dimension, producing a fixed-size representation regardless of $n$:
\begin{equation}
\mathbf{e}_{uv}^{\mathrm{conv}} = \mathrm{GAP}(\mathbf{H}_5) = \frac{1}{n}\sum_{t=1}^{n} \mathbf{H}_5[\,:\,,t\,] \in \mathbb{R}^{64}.
\label{eq:gap}
\end{equation}

\subsection{Edge Type Embedding and Merge}

Each directed edge $(u, v)$ is assigned a discrete edge type $\tau_{uv}$ from a 7-class taxonomy based on the positions of $u$ and $v$ relative to $X$ and $Y$. A learnable embedding table maps each type to a vector $\mathbf{t}_{\tau_{uv}} \in \mathbb{R}^{64}$. The final edge representation combines the Conv1D embedding with the type embedding via addition, which preserves the representation dimension:
\begin{equation}
\tilde{\mathbf{e}}_{uv} = \mathbf{e}_{uv}^{\mathrm{conv}} + \mathbf{t}_{\tau_{uv}}.
\label{eq:merge}
\end{equation}

\subsection{Binary Edge Head}

A linear layer projects each $\tilde{\mathbf{e}}_{uv}$ to 2 logits corresponding to the binary label of whether the directed edge $u \to v$ exists in the ground-truth DAG. This edge head is used only during training as an auxiliary supervision signal; its predictions are discarded at inference. The role of this auxiliary supervision is discussed in Section~\ref{sec:loss}.

% =====================================================================
\section{Structural Attention Bias}
\label{sec:struct_bias}
% =====================================================================

This section describes the second contribution of this thesis: a topology-aware attention mechanism that injects causal structural priors into self-attention with only six learnable scalars per head.

\subsection{Motivation}

Standard multi-head self-attention assigns attention weights between all pairs of edge representations purely from learned query-key dot products, with no consideration of their topological relationship in the causal graph. This ignores a strong prior: edges that share an endpoint or form a directed chain through a node of interest are structurally more relevant to each other than unrelated background edges. For example, in the chain $X \to K \to Y$, the edges $X \to K$ and $K \to Y$ should attend to each other more strongly than to a background edge $A \to B$.

\subsection{Six Topological Relationship Types}

For each ordered pair of edges $(i, j)$ in the graph, a topological relationship type $\tau(i,j) \in \{0,\ldots,5\}$ is assigned based on shared endpoints. Given edges $e_i = (u_1, v_1)$ and $e_j = (u_2, v_2)$, the six types are defined in Table~\ref{tab:struct_types}. Each type carries a distinct causal interpretation: reverse pairs represent the same variable pair in opposite directions (causal asymmetry signal); shared source pairs indicate a common cause structure (fork); shared target pairs indicate a common effect structure (collider); forward and backward chain pairs indicate mediator pathways; unrelated pairs carry no topological structure.

\begin{table}[ht]
	\centering
	\caption{Six topological relationship types between edge pairs in the structural attention bias. Each type is determined by whether the two edges share a source or target node, encoding a distinct causal structural relationship. The learned scalar bias $b_\tau$ allows the model to selectively strengthen or weaken attention between edge pairs according to their causal topology.}
	\label{tab:struct_types}
	\begin{tabular}{lll}
		\toprule
		\textbf{Type} & \textbf{Condition} & \textbf{Causal interpretation} \\
		\midrule
		Reverse        & $u_1 = v_2,\; v_1 = u_2$ & Same pair, opposite direction \\
		Shared source  & $u_1 = u_2,\; v_1 \neq v_2$ & Common cause (fork) structure \\
		Shared target  & $v_1 = v_2,\; u_1 \neq u_2$ & Common effect (collider) structure \\
		Forward chain  & $v_1 = u_2$ & Mediator chain $u_1 \to v_1 \to v_2$ \\
		Backward chain & $v_2 = u_1$ & Reverse mediator chain \\
		Unrelated     & otherwise & No shared endpoint \\
		\bottomrule
	\end{tabular}
\end{table}

\subsection{Learned Scalar Bias per Topological Type}

The attention logit between edges $i$ and $j$ is augmented with a learnable scalar bias specific to their topological relationship:
\begin{equation}
A_{ij} = \frac{\mathbf{q}_i \cdot \mathbf{k}_j + b_{\tau(i,j)}}{\sqrt{d}},
\label{eq:struct_bias}
\end{equation}
where $\mathbf{q}_i, \mathbf{k}_j \in \mathbb{R}^d$ are the standard query and key projections, $d$ is the head dimension, and $b_{\tau} \in \mathbb{R}$ is a single learnable scalar for each of the six types. This introduces only 6 additional parameters per attention head while encoding explicit causal structural prior knowledge. All bias scalars are initialized to zero, so training begins from standard unbiased self-attention, preserving learned representations from early epochs while gradually adjusting attention weights according to causal topology. Two layers of multi-head self-attention with structural bias are applied sequentially to aggregate information across the full edge set.

\subsection{Why Structured Bias Outperforms Unconstrained Capacity}

Several other architectural extensions -- parallel observation cross-attention branches, observation-level transformers, node-centric aggregation -- all produce results at or below the 3-channel baseline. The distinguishing factor is that structural attention bias adds a structured inductive prior with a minimal parameter budget (6 scalars per head), rather than additional unconstrained capacity. The bias does not learn a new representation from scratch -- it only learns the relative priority among topologically distinct edge pairs. This is validated by the experimental result: structural bias improves Balanced Accuracy by +2.65\% over the same input channels without it, a gain larger than any other single architectural change in this thesis.

% =====================================================================
\section{Node Visual Representation}
\label{sec:node_pipeline}
% =====================================================================

This section describes the third contribution of this thesis: a 2D scatter density pipeline that provides complementary visual features to the 1D edge tensor representation.

\subsection{Motivation}

The 1D edge tensor representation is inherently pairwise: it describes the marginal relationship between $u$ and $v$ rather than the joint visual pattern of a node $K$ with respect to both $X$ and $Y$ simultaneously. The node visual representation complements the 1D pipeline by rendering the joint scatter distribution of $K$ against $X$ and $Y$ as a two-dimensional density image. This captures visual spatial structure -- clustering, curvature, heteroscedasticity, and residual patterns -- that is invisible in any single 1D marginal curve.

\subsection{8-Channel Scatter Density Image}

For each non-$X$/$Y$ node $K$, an 8-channel $32 \times 32$ image $\mathbf{I}_K \in \mathbb{R}^{8 \times 32 \times 32}$ is constructed from the $n = 1000$ observations. The 8 channels are grouped into two sets with separate smoothing parameters.

Channels 1--4 render the raw scatter density with $\sigma_{\mathrm{raw}} = 4.0$: $\mathbf{I}_K^{(1)} = \rho(K, X)$, $\mathbf{I}_K^{(2)} = \rho(K, Y)$, $\mathbf{I}_K^{(3)} = \rho(X, K)$, $\mathbf{I}_K^{(4)} = \rho(Y, K)$, where $\rho(a, b)$ denotes the smoothed $32 \times 32$ joint density image of variables $a$ (horizontal) and $b$ (vertical), each axis min-max normalized to $[-1, 1]$. The four views provide the full bidirectional visual context of $K$ with respect to $X$ and $Y$. A larger smoothing $\sigma_{\mathrm{raw}} = 4.0$ is used to smooth out observation noise while preserving the coarse spatial structure of the joint distribution.

Channels 5--8 encode directional asymmetry with $\sigma_{\mathrm{ANM}} = 2.0$: $\mathbf{I}_K^{(5)} = \rho(K, r_{X \to K})$, $\mathbf{I}_K^{(6)} = \rho(X, r_{K \to X})$, $\mathbf{I}_K^{(7)} = \rho(K, r_{Y \to K})$, $\mathbf{I}_K^{(8)} = \rho(Y, r_{K \to Y})$, where $r_{u \to v}$ denotes the ANM residual of $v$ regressed on $u$. If the true causal direction is $u \to v$, the residual $r_{u \to v}$ is approximately independent of $u$, producing a uniform horizontal scatter; the reverse-direction residual retains visible nonlinear structure. A finer smoothing $\sigma_{\mathrm{ANM}} = 2.0$ is used to preserve the fine-grained directional asymmetry patterns that coarser smoothing would obscure.

An earlier variant with 12 channels (adding kernel coefficient maps) scored 80.42\% versus 80.47\% for the 8-channel version. The kernel coefficient represents a derived linear relationship with no meaningful spatial structure along the observation axis, making Conv2D's spatial inductive bias inappropriate for these channels.

\subsection{Hierarchical Conv2D Feature Extractor}

The 8-channel $32 \times 32$ image $\mathbf{I}_K$ is processed by a hierarchical Conv2D network that progressively downsamples the spatial resolution: $32 \times 32 \to 16 \times 16 \to 8 \times 8 \to 4 \times 4 \to \mathbf{e}_K^{\mathrm{2D}} \in \mathbb{R}^{64}$. Each Conv2D block uses Group Normalization and GELU activations consistent with the 1D pipeline. The final AdaptiveAvgPool collapses the $4 \times 4$ spatial map to a $d = 64$ dimensional node visual embedding.

% =====================================================================
\section{Fusion and Node Classification}
\label{sec:fusion}
% =====================================================================

This section describes how the edge context and visual representations are combined for the final 8-class prediction.

\subsection{Combining Edge Context with Visual Representation}

For each node $K$, the final classification embedding is formed by combining five sources of information: the four directional edge context embeddings from the Edge Context Encoder and the visual embedding from the Node Visual Classifier. Let $\mathbf{e}_{K \to X}$, $\mathbf{e}_{K \to Y}$, $\mathbf{e}_{X \to K}$, $\mathbf{e}_{Y \to K} \in \mathbb{R}^{64}$ be the four attended edge embeddings, and let $\mathbf{e}_K^{\mathrm{2D}} \in \mathbb{R}^{64}$ be the node's visual embedding. These are combined via a learned linear merge operator:
\begin{equation}
\mathbf{h}_K = \mathrm{MergeOperator}\!\left(\mathbf{e}_{K \to X},\; \mathbf{e}_{K \to Y},\; \mathbf{e}_{X \to K},\; \mathbf{e}_{Y \to K},\; \mathbf{e}_K^{\mathrm{2D}}\right) \in \mathbb{R}^{64}.
\label{eq:node_merge}
\end{equation}
The node head then maps $\mathbf{h}_K$ to 8-class logits:
\begin{equation}
\hat{\mathbf{y}}_K = \mathrm{Linear}_{64 \to 8}(\mathbf{h}_K).
\label{eq:node_head}
\end{equation}

\subsection{Necessity of Edge Context}

Table~\ref{tab:edge_context_necessity} presents ablation results confirming that the edge context is structurally necessary, not merely supplementary. When the entire 1D edge pipeline is removed and nodes are classified from their 2D visual embeddings alone, Balanced Accuracy collapses to 51.51\% -- a drop of nearly 29 percentage points from the full model's 80.47\%. A variant using mean pooling of all node embeddings as a substitute for edge context (v28b) achieves only 46.39\%. This confirms that the 2D visual representation captures local distributional structure of $K$ relative to $X$ and $Y$, but cannot reason about $K$'s causal role within the larger graph -- a task that requires the global relational information encoded by the attended edge embeddings.

\begin{table}[ht]
	\centering
	\caption{Ablation study on the necessity of edge context for node classification. The full model (v26e) combines 1D edge context with 2D visual features, achieving 80.47\% Balanced Accuracy. Removing the 1D edge pipeline and using 2D visual features alone (v26c) collapses performance to 51.51\%, a 28.96-point drop. Replacing edge context with mean-pooled node embeddings (v28b) further degrades to 46.39\%. The binary edge loss, by supervising the edge pipeline with ground-truth adjacency matrix entries, ensures these embeddings remain causally well-structured throughout training.}
	\label{tab:edge_context_necessity}
	\begin{tabular}{lcc}
		\toprule
		\textbf{Variant}           & \textbf{Balanced Accuracy} & \textbf{Gap to Full} \\
		\midrule
		v26e (full model)         & \textbf{80.47\%}           & --                   \\
		v28b (2D only + mean pool) & 46.39\%                     & $-$34.08\%           \\
		v26c (2D only)             & 51.51\%                     & $-$28.96\%           \\
		\bottomrule
	\end{tabular}
\end{table}

% =====================================================================
\section{Dual-Head Loss Weighting}
\label{sec:loss}
% =====================================================================

This section describes the fourth contribution of this thesis: a dual-head loss weighting scheme that ensures the edge context encoder remains causally well-supervised throughout training.

\subsection{Dual-Head Loss with Lambda Weighting}

The model is trained with two loss terms. The primary node loss $\mathcal{L}_{\mathrm{node}}$ is a weighted cross-entropy over the 8 causal role classes, with class weights set inversely proportional to class frequency in the training set:
\begin{equation}
\mathcal{L}_{\mathrm{node}} = -\sum_{c=1}^{8} w_c \cdot \mathbf{1}[y = c] \cdot \log \hat{p}_c, \qquad w_c \propto \frac{1}{n_c},
\label{eq:node_loss}
\end{equation}
where $n_c$ is the number of training samples of class $c$, $y$ is the ground-truth label, and $\hat{p}_c$ is the predicted probability for class $c$. The auxiliary edge loss $\mathcal{L}_{\mathrm{edge}}$ is a binary weighted cross-entropy over all $p(p-1)$ directed edges, predicting whether each edge exists in the ground-truth DAG. The total training loss is:
\begin{equation}
\mathcal{L} = \lambda \cdot \mathcal{L}_{\mathrm{edge}} + (1 - \lambda) \cdot \mathcal{L}_{\mathrm{node}}, \quad \lambda = 0.7.
\label{eq:total_loss}
\end{equation}

\subsection{Rationale for $\lambda = 0.7$}

The choice of $\lambda = 0.7$ assigns the majority of the training gradient to the auxiliary edge loss. This is counterintuitive -- the edge head is discarded at inference and the primary objective is node classification -- but it is empirically necessary. The 2D node classifier receives its structural context entirely from the attended edge embeddings. If $\lambda$ is small (node loss dominates early training), the edge pipeline receives insufficient gradient and converges to representations optimized for the node head's immediate needs rather than for causally meaningful edge structure, producing degenerate context that degrades the node merger quality. Increasing $\lambda$ to 0.7 ensures the edge embeddings are well-supervised throughout training. This mechanism is confirmed by ablation: a configuration favoring node loss ($\lambda = 0.3$) achieves 78.77\%, while $\lambda = 0.7$ achieves 80.47\% -- a difference of +1.70\% with identical input channels and architecture, differing only in the loss weighting.

% =====================================================================
\section{X/Y Remap Data Augmentation}
\label{sec:xyaug}
% =====================================================================

This section describes the fifth contribution of this thesis: a data augmentation strategy that multiplies the effective training set by approximately 11x without additional data collection.

\subsection{Strategy}

The training set contains 23,500 dataset instances, each with a designated pair $(X, Y)$ satisfying $X \to Y$ in the ground-truth DAG. However, a dataset instance with $p$ variables contains multiple valid directed edges in the ground-truth DAG, each of which could serve as a treatment-outcome pair. The X/Y remap augmentation exploits this symmetry: for each directed edge $A \to B$ in the ground-truth DAG, a new training sample is generated by treating $(A, B)$ as the new $(X', Y')$ and recomputing 8-class labels for all remaining nodes $K'$ according to their causal relationship with $A$ and $B$. This produces approximately 11.2x more training samples, yielding a training set of approximately 263,000 samples. The augmentation is justified by the symmetry of the classification problem: the 8 causal roles are defined purely by the local structure of the graph around the designated $(X, Y)$ pair, and this local structure occurs with equal validity for any directed edge in the DAG.

\subsection{Efficient Implementation}

Naively recomputing the 8-channel edge tensors for each remap would multiply the preprocessing cost by approximately 11x. The key efficiency insight is that the kernel regression computation (channels 3--8) depends only on the observation matrix $\mathbf{D}$, not on the choice of $(X, Y)$. Therefore, coefficient maps and residual maps are computed once per graph. For each remap, only three operations are required: edge type re-assignment (cheap name-substitution), node image reconstruction from cached maps (cheap histogram operations), and label re-derivation from the adjacency matrix (a lookup table query). No kernel regression is recomputed.

Gains from X/Y remap augmentation are configuration-dependent. On v8b (3-channel baseline), augmentation provides +3.32\% improvement. On v11 (8-channel, structural attention), augmentation provides +1.55\%. On v26e (full dual-pipeline), augmentation provides +0.72\%. The diminishing returns reflect that richer base representations benefit less from additional data multiplication, though the augmentation remains beneficial across all configurations.

% =====================================================================
\section{Training Configuration}
\label{sec:training_config}
% =====================================================================

All models are trained with the following fixed configuration: AdamW optimizer~\cite{loshchilov2019adamw} with cosine annealing learning rate schedule~\cite{loshchilov2017sgdr}, initial learning rate $10^{-3}$, batch size 64, 10 training epochs on the augmented 263K dataset (equivalent to approximately 110 epochs over the base 23,500 samples), and 32-bit floating-point precision throughout. A Gaussian noise augmentation of standard deviation 0.01 is applied to edge tensors during training to improve generalization. Training uses Distributed Data Parallel across two NVIDIA RTX 5880 GPUs. The model has approximately 280K parameters in total (edge encoder: approximately 188K; node Conv2D extractor: approximately 92K; node merge and head: negligible).

% =====================================================================
\section*{Chapter Summary}
\addcontentsline{toc}{section}{Chapter Summary}
% =====================================================================

This chapter presented the five contributions that extend the first-place deep learning baseline from 76.70\% to 81.082\% Balanced Accuracy. The first contribution, an 8-channel edge tensor, extends the 3-channel representation with multi-bandwidth kernel regression and ANM residuals, preserving the full shape of conditional relationships that scalar features irreversibly discard. The second contribution, a structural attention bias, injects topology-aware causal priors into self-attention with only six learned scalars, encoding chain, fork, and collider relationships directly and outperforming failed experiments that added thousands of unconstrained parameters. The third contribution, a 2D scatter density pipeline, complements the 1D edge tensors by rendering joint distributions as visual images that capture multimodality, heteroscedasticity, and residual patterns invisible to pairwise analysis. The fourth contribution, a dual-head loss weighting scheme with $\lambda = 0.7$, ensures the edge context encoder remains causally well-supervised throughout training, producing richer structural embeddings that the node classifier depends on. The fifth contribution, X/Y remap data augmentation, multiplies the training set by approximately 11x without additional data collection, with the largest gains on difficult minority classes. The combined system achieves 81.082\% Balanced Accuracy on the CrunchDAO private test set.